\chapter{Integration by Parts}\label{ch:parts}

Integration by parts is the product rule integrated, and the chapter is in three movements:
the derivation and the heuristic for using it; the engine's finder, whose every candidate was
correct and every candidate refused; and what the refusal taught about normal forms, which
became the shape of the checker's theorem.

\section{The formula}

\begin{theorem}
  If $u, v$ are differentiable with continuous derivatives, then
  $\int u\,v'\,dx = uv - \int u'\,v\,dx$.
\end{theorem}
\begin{proof}
  $(uv)' = u'v + uv'$ by the product rule, so $uv$ is an antiderivative of $u'v + uv'$, and by
  linearity $\int uv' = uv - \int u'v$.
\end{proof}

Written $\int u\,dv = uv - \int v\,du$, the formula trades one integral for another, and the art
is choosing $u$ so that the new integral is easier. The heuristic LIATE --- logarithm, inverse
trigonometric, algebraic, trigonometric, exponential --- ranks candidates for $u$: take as $u$
the factor earliest in the list. The reason is that differentiation \emph{simplifies} a
logarithm (to $1/x$) and a polynomial (lowers its degree), while it leaves an exponential or a
sine as it was; so $u$ should be the factor that gets simpler, and $dv$ the one that is easy to
integrate.

\begin{example}
  $\int x e^x\,dx$: $u = x$, $dv = e^x dx$, so $du = dx$, $v = e^x$, and the integral is
  $x e^x - \int e^x\,dx = x e^x - e^x$.
\end{example}

\begin{example}
  $\int \ln x\,dx$: $u = \ln x$, $dv = dx$, so $du = dx/x$, $v = x$, and the integral is
  $x \ln x - \int 1\,dx = x \ln x - x$. This is the table entry of Chapter~\ref{ch:anti}.
\end{example}

\begin{example}
  $\int x^2 e^x\,dx$ needs two rounds: $x^2 e^x - 2\int x e^x\,dx = x^2 e^x - 2(x e^x - e^x)$.
  Each round lowers the degree by one; a polynomial of degree $n$ needs $n$ rounds.
\end{example}

\begin{example}[The solve-for trick]
  $\int e^x \sin x\,dx$: two rounds return to the original integral with a sign,
  $\int e^x \sin x = e^x \sin x - e^x \cos x - \int e^x \sin x$, and \emph{solving} gives
  $\frac{1}{2} e^x(\sin x - \cos x)$. This is the one shape of by-parts the engine does not
  attempt, because its finder never solves an equation.
\end{example}

\section{The finder}

\begin{minted}{lean}
-- integration by parts, a logarithm first, else a power of the variable
if fuel = 0 then none else
let isLog : Expr → Bool | .fn "ln" [_] => true | _ => false
let isPoly : Expr → Bool | .var y => y == x | .pow (.var y) (.num n) => y == x && n.isInt && !n.isNeg | _ => false
let (u, i) ← (es.zipIdx.find? (isLog ·.1)) <|> (es.zipIdx.find? (isPoly ·.1))
let dv := without es i
let (v, vsteps) ← anti simp x fuel dv
let du ← simp (D u x)
let vdu ← simp (.mul [v, du])
let (inner, isteps) ← anti simp x (fuel - 1) vdu
let F := Expr.sub (.mul [u, v]) inner
return (F, #[step "int.by-parts" "…" F] ++ vsteps ++ isteps)
\end{minted}

This is LIATE with two letters --- a logarithm first, else a power of the variable --- and a
depth of three rounds, enough for a cubic times an exponential. The recursion on
\lc{fuel - 1} is the one place in the finder with a budget; it bounds the depth of a
\emph{search}, not the length of an evaluation, and the search failing is a refusal.

\section{Every candidate refused}

When by-parts was first added, every candidate it produced was correct and every one was
refused by the checker. The reason is instructive. The candidate for $\int x e^x$ is
$x e^x - e^x$; its derivative normalizes to $x e^x + e^x - e^x$, and then --- since the simplifier
collects like terms --- to $x e^x$. That one passes. But for $\int x \sin x$ the candidate is
$-x \cos x - \int (-\cos x)\,dx = -x\cos x + \sin x$, and its derivative is
$-\cos x + x \sin x + \cos x$, which the simplifier does collect. The failures came from the
minus sign in front of a sum: $-(a + b) + b$ is a \emph{normal form} of the pipeline, because
the pipeline never distributes a numeral over a sum. It cannot: distribution grows terms, and
the termination ordering of Chapter~\ref{ch:wf} forbids it. Distribution is the job of
\texttt{expand}, which is a function, not a rule (Chapter~\ref{ch:poly}).

So the checker was changed, in the only direction that keeps it sound: both the derivative of
the candidate and the integrand are expanded with \lc{Expand.dist} before they are normalized
and compared. \lc{dist} is proved value-preserving, for the derivative semantics as well as the
plain one:

\begin{minted}{lean}
theorem dist_soundD : ∀ (e : Expr) (ρ : EnvR), evalD ρ (Expand.dist e) = evalD ρ e
\end{minted}

This is the second instance of the generic proof of Chapter~\ref{ch:poly} --- the structure
\lc{Sem} instantiated with \lc{evalD} --- and it is why the theorem \lc{cmdIntegrate_spec}
compares \lc{norm (Expand.dist g)} with \lc{norm (Expand.dist f)} rather than \lc{g} with
\lc{f}. Nothing was weakened. A candidate is still accepted only if its derivative \emph{is} the
integrand; the identification is just allowed to use the distributive law, which is a theorem.

\section{Powers of sine and cosine}

$\int \sin^3 x\,dx$ is by parts with the solve-for trick built in. Take $u = \sin^{m-1} x$ and
$dv = \sin x\,dx$:
\begin{align*}
  \int \sin^m x\,dx &= -\sin^{m-1} x \cos x + (m-1)\int \sin^{m-2} x \cos^2 x\,dx \\
    &= -\sin^{m-1} x \cos x + (m-1)\int \sin^{m-2} x\,dx - (m-1)\int \sin^m x\,dx,
\end{align*}
using $\cos^2 = 1 - \sin^2$, and solving for the integral gives the \emph{reduction formula}
\[
  \int \sin^m x\,dx = -\frac{\sin^{m-1} x \cos x}{m} + \frac{m-1}{m} \int \sin^{m-2} x\,dx .
\]
The same computation with a factor $\cos^n x$ carried along gives
\[
  \int \sin^m x \cos^n x\,dx = -\frac{\sin^{m-1} x\cos^{n+1} x}{m+n} + \frac{m-1}{m+n}\int \sin^{m-2} x\cos^n x\,dx,
\]
and its mirror reduces the cosine. Two steps take $\sin^3$ to $\sin$; the engine's finder applies
exactly these formulas (\rn{int.trig-power}), each recorded as a step, until only first powers
remain for the table and substitution to finish.

The check has a new difficulty here. Differentiate the answer for $\int \sin^2 x\,dx$, namely
$-\frac{1}{2}\sin x \cos x + \frac{1}{2} x$: the derivative is $\frac{1}{2} - \frac{1}{2}\cos^2 x +
\frac{1}{2}\sin^2 x$, which equals $\sin^2 x$ --- but only \emph{modulo} $\sin^2 + \cos^2 = 1$, an
identity the simplifier does not apply, because neither direction of it decreases the ordering.
So the comparison first rewrites every $\cos^k u$ with $k \ge 2$ as $(1 - \sin^2 u)^{k/2}$ times
$\cos u$ if $k$ is odd, on both sides, then expands and simplifies. After that rewrite a
trigonometric polynomial has cosine to at most the first power, which is a normal form for the ring
$\mathbb{R}[s, c]/(s^2 + c^2 - 1)$, and two equal trigonometric polynomials become the same term.
The same pass rewrites $(e^u)^k$ as $e^{ku}$, which is why $\int (e^x)^2\,dx$ verifies. The pass
is a total function like \lc{dist}, proved value-preserving in both semantics
(\lc{identNorm_sound}, \lc{identNorm_soundD}) from Mathlib's \lc{Real.cos_sq'} and the
definition of a real power of a positive base; the checker's theorem names it alongside
\lc{dist}. Nothing was weakened: a candidate is still accepted only if its derivative
\emph{is} the integrand, with the Pythagorean identity and distribution allowed in the
identification because they are theorems.

\section{The tangent}

A second refusal became a rule. $\int \tan x$ was refused because the candidate $-\ln \cos x$
has derivative $\sin x \cdot (\cos x)^{-1}$, and the engine did not know that this is $\tan x$.
The right fix was not to teach the checker to accept it --- that would weaken the check --- but to
add the identity $\sin u / \cos u = \tan u$ to the simplifier, prove it terminates under both of
the engine's orderings, and prove it sound over $\mathbb{R}$ (Chapter~\ref{ch:explog}: it is
unconditional because Mathlib's tangent is defined as the quotient). Three proofs for one
identity is the price of a rule, and it is why rules are added when a checker refuses something
true and not before.

\section{What is still refused}

$e^x \sin x$ needs the solve-for trick for a product of two functions that never simplify.
$\frac{1}{x^2 + 1}$ needs an arctangent, which the engine does not have. Each is refused with the candidate it tried and the derivative it got, and none is
a wrong answer. The finder can be extended without touching a proof; the checker's theorem
covers whatever it produces.

\begin{trybox}
\begin{minted}{text}
integrate(x * exp(x), x)
integrate(x^2 * exp(x), x)
integrate(ln(x), x)
integrate(sin^3(x), x)
integrate(x * sin^3(x), x)
integrate(tan(x), x)
integrate(exp(x) * sin(x), x)
\end{minted}
  The first three succeed by parts; open the compare step to see both sides expanded. The
  fourth and fifth use the reduction formula, the fifth inside by parts. The sixth is the
  tangent. The last is refused, with the reason.
\end{trybox}

\section{Exercises}

\begin{exercisebox}[Tabular integration]
  Show that for a polynomial $p$ of degree $n$, $\int p(x) e^x\,dx = e^x \sum_{k=0}^n (-1)^k
  p^{(k)}(x)$, by induction on $n$ using by parts.
\end{exercisebox}

\begin{exercisebox}[The solve-for trick as a rule]
  Design a finder step for $\int e^{ax} \sin(bx)\,dx$ that guesses $e^{ax}(A \sin bx + B\cos bx)$
  and solves for $A, B$ by differentiating. What would its status be, and why does the checker's
  theorem need no change?
\end{exercisebox}

\begin{exercisebox}[Normal forms]
  Give a term that the pipeline leaves as $-(a + b) + b$ and show, step by step, what
  \lc{Expand.dist} followed by the simplifier does to it. Which rule of Chapter~\ref{ch:poly}
  fires last?
\end{exercisebox}
